\DeclareFontFamily{U}{matha}{\hyphenchar\font45}
\DeclareFontShape{U}{matha}{m}{n}{
    <5> <6> <7> <8> <9> <10> gen * matha
    <10.95> matha10 <12> <14.4> <17.28> <20.74> <24.88> matha12
}{}
\newcommand\bbl{\mathrel{\ooalign{\usefont{U}{matha}{m}{n}\symbol{\string"CE}}}}

We continue the development of sized types with parametric quantifiers over sizes.
In previous work \cite{licentiate, parametric-quantifier}, the type of sizes was interpreted in the semantics as the natural numbers.
This means that the following axiom for small types was restricted to the case where \(A\) is finite: \begin{align*}
    \Pi{x : A} \. \exists i . B(x, i)&\simeq\exists i\. \Pi {x : A} \. \exists j < i . B(x, j). \tag{\(\exists\) parametric over \(\Pi\)}
\end{align*}
Our main contribution is a model where the type of sizes is instead interpreted as a larger ordinal.
Indeed, for infinite \(A\), the left-to-right direction of the equivalence intuitively shows that limit ordinals exist.
This model allows us to remove the restriction that \(A\) is finite, which enables us to define all (co)inductive types, not just the finitely-branching ones.

Our model is based on realisability instead of the reflexive graphs used in previous work \cite{licentiate, parametric-quantifier}.
This semantics fits well with the intuition of parametric quantifiers: \(\exists i.\,A(i)\) behaves like the union of the \(A(i)\), while \(\forall i.\,A(i)\) behaves like the intersection of the \(A(i)\).
The model also shows that our syntax is consistent with two additional principles: an impredicative universe of sets, and uniqueness of identity proofs.
However, we do not use these properties in the syntax, which makes our construction of (co)inductive types more widely applicable and compatible with homotopy type theory.

\subparagraph*{Discussion and Future work.}

A main motivation for our work is obtaining a consistent implementation of sized types in proof assistants.
Encoding (co)inductive types using quantifiers would allow the closure ordinal \(\infty\) to be removed from \texttt{Agda}, which is a source of inconsistency of the current implementation.
However, proof assistants should satisfy desirable computational properties like decidability and canonicity, and this requires computational behaviour for our parametricity axioms.
So, to use our results for this purpose, we need to connect them with notions of internal parametricity that compute, such as those based on bridge types \cite{parametric-quantifier}, which have also been developed in combination with cubical type theory \cite{cavallo2021internal,parametricity-agda}.

In addition, we need a computational version of the $\beta$-rule for the fixpoint operator.
We have only assumed a propositional version of this rule because a na\"ive definitional version of this rule breaks normalisation due to unrestricted unfolding.
For example, it would be possible to disallow unfolding under binders, only allowing unfolding when a smaller size exists in context, or to ignore computation on sizes entirely.
However, each of these options come with drawbacks and it is not clear to us what the most principled approach is.

Because we include suprema sizes, it is not generally possible to keep inequality decidable.
In the current \(\sf{Agda}\) implementation, \(i<j\) is a judgement that is automatically checked when applying a function \(f : (i : \sf{Size}_{<j}) \to A\). However, without decidability of inequality we need to ask the user for a term of the type \(i<j\), which increases the amount of manual work.
One partial solution is to have two versions of inequality: a propositional type \(i<j\) and a definitional judgemental fragment \(i\bbl j\) (similar to the two versions of equality, \(=\) and \(\equiv\)).
Then we could add a term \(\sf{ineq}\) so that \(\sf{ineq}:i<j\) holds whenever \(i\bbl j\) (similar to the term \(\sf{refl}\) where \(\refl:a=_Aa'\) holds whenever \(a\equiv a'\)).
This would allow the user to fill in a simple term when the computer can decide the inequality, while still allowing more complicated terms when using suprema sizes.

On the semantic side, it is worth investigating whether we can generalise our realisability model.
We restricted the model to the pca of partial recursive functions, as it allowed us to interpret the type of sizes as the first uncountable ordinal.
In principle, it seems that the model could generalise to other pcas, as long as the type of sizes is interpreted as a sufficiently large ordinal: one of larger cardinality than the pca. The realizability model could also be constructed using more general notions of computation, such as the monadic combinatory algebras introduced in \cite{Cohen-et-al-2025}, to allow for an integration of sizes with effects.

Another question is whether we can model both sizes and the univalence axiom. A possibly suitable model for this is given by cubical assemblies \cite{cubical-assemblies}, which form a model of a type theory with a univalent and impredicative universe. 
It seems likely that this construction can be extended to include sizes and parametric quantifiers in a similar way as in this work, but a more careful investigation is needed to confirm this.

We would also like to investigate how realisability models can model internal parametric quantifiers and related modalities such as irrelevance in dependent type theory more generally.
The inclusion from sets to assemblies assigns a set the complete realisability relation, which intuitively makes elements computationally indistinguishable and enables parametricity for sizes.
It seems reasonable that this construction would work for a more general notion of parametric quantifier in the theory.